\subsection*{Key Findings}
\begin{itemize}
\item Orbital Mapper shows high structural complexity with low imputation dependence.
\item Thermal Mapper shows high complexity but high imputation dependence.
\item Adding derived density slightly modifies, and in this run reduces, Mapper complexity.
\end{itemize}

\subsection*{Promising Signals}
\begin{itemize}
\item The orbital PCA Mapper is a high-priority graph for scientific inspection: it combines nontrivial cycle structure with low imputation dependence.
\item Highlighted node catalog available with 495 node-level interpretation records.
\item Connected-component summaries available for 6 principal graphs.
\end{itemize}

\subsection*{Risky Configurations}
\begin{itemize}
\item The thermal Mapper should be treated cautiously because more than half of its nodes exceed the high-imputation threshold.
\end{itemize}

\subsection*{Density Effect}
Adding derived density reduces cycle complexity in the joint space, suggesting that density acts as a regularizing coordinate rather than introducing additional branching.

\subsection*{Lens Effect}
Results are lens-sensitive; pca2 should remain the primary interpretation layer, while density/domain are sensitivity probes.

\subsection*{Next Steps}
\begin{itemize}
\item Bootstrap stability for principal graphs.
\item Null models based on column-wise shuffling.
\item Astrophysical review of highlighted nodes and components.
\item Bootstrap was not run in this execution.
\item Null models were not run in this execution.
\item Multi-method Mapper comparison was skipped or partial because not all imputation inputs were available.
\end{itemize}
